\documentclass{report}
\usepackage{amsmath,amssymb}
\setlength{\parindent}{0mm}
\setlength{\parskip}{1em}
\begin{document}
\begin{center}
\rule{6in}{1pt} \
{\large Jon Calhoun \\
{\bf How Silent Errors Propagate in Iterative Methods}}

201 N Goodwin Ave \\ Urbana \\ IL 61801
\\
{\tt jccalho2@illinois.edu}\\
Luke Olson\\
Marc Snir\end{center}


As simulation size and machine complexity grows on next generation HPC
systems, machine errors are expected to increase. In particular, silent
data corruption (SDC), which occurs due to cosmic radiation striking
hardware components causing the state of a transistor to flip, remains a
concern key concern. SDC in iterative methods can lead to extra
iterations required to find a solution, can convergence to an incorrect
solution, or cause an application crash. Understanding how SDC propagates
through iterative methods can lead to better mitigation systems that in
turn increases effective system utilization.

In this talk, we highlight SDC in several iterative computations over a
range of applications. For example we consider miniapps, including CoMD,
as well as benchmarks from NAS and the graph500 suites. From this we show
how SDC propagates through different operations, suggesting different
algorithm-based mitigation strategies to help reduce the impact of SDC.


\end{document}
